%------------------------------------------------------------------------------%
\documentclass[fleqn,utf8,aspectratio=169,12pt]{beamer}

\usepackage{integracao_e_series}

\title{Integrais de Funções Trigonométricas}
\subtitle{1 -- Senos e Cossenos}

%------------------------------------------------------------------------------%
\begin{document}

\addtitle

%------------------------------------------------------------------------------%
\section{Integrais de Funções Trigonométricas}

\begin{frame}{Objetivo}
  Estratégias para calcular
  \[
    \int \sin^m(x) \cos^n(x)\,dx
  \]

  \pause
  \[
    \int \tan^m(x) \sec^n(x)\,dx
  \]

  \pause
  \[
    \int \cot^m(x) \csc^m(x)\,dx
  \]
\end{frame}

%------------------------------------------------------------------------------%
\begin{frame}{Substituição Simples -- Abordagem Ingenua}
\begin{align*}
  \onslide<+->{I & = \int \sin^3(x) \cos^2(x)dx                     &  & \cos^2(x) = 1 - \sin^2(x)                      \\[2mm]}
  \onslide<+->{  & = \int \sin^3(x) \left( 1 - \sin^2(x) \right) dx &  &                                                \\[2mm]}
  \onslide<+->{  & = \int \sin^3(x) - \sin^5(x)\,dx                 &  &                                                \\[2mm]}
  \onslide<+->{  & = \int \sin^3(x)\,dx - \int\sin^5(x)\,dx         &  & u = \sin(x) \qquad du = {\color{red}\cos(x)}dx \\[2mm]}
  \onslide<+->{  & = \;?                                            &  &}
\end{align*}
\end{frame}

%------------------------------------------------------------------------------%
\begin{frame}{Relações Trigonometrias Úteis}
\vspace{\baselineskip}
\begin{columns}
\column{0.6\linewidth}
  \onslide<+->{Relação fundamental
  \vspace{-0.3\baselineskip}
  \[
    \sin^2(x) + \cos^2(x) = 1
  \]}
  \onslide<+->{Decorrências da relação fundamental
  \vspace{-0.3\baselineskip}
  \begin{gather*}
    \tan^2(x)+ 1  = \sec^2(x) \\
    1 + \cot^2(x) = \csc^2(x)
  \end{gather*}}
  \onslide<+->{Seno e cosseno da soma
  \vspace{-0.3\baselineskip}
  \begin{gather*}
    \sin(x + y) = \sin(x)\cos(y) + \cos(x)\sin(y) \\
    \cos(x + y) = \cos(x)\cos(y) - \sin(x)\sin(y)
  \end{gather*}}
\column{0.4\linewidth}
  \onslide<+->{Arco duplo
  \begin{gather*}
    \sin(2x) = 2\sin(x) \cos(x) \\
    \cos(2x) = \cos^2(x)-\sin^2(x)
  \end{gather*}}

  \onslide<+->{Arco metade
  \begin{gather*}
    \sin^2(x) = \frac{1}{2}\big(1-\cos(2x)\big) \\[2mm]
    \cos^2(x) = \frac{1}{2}\big(1+\cos(2x)\big)
  \end{gather*}}
\end{columns}
\end{frame}

%------------------------------------------------------------------------------%
\begin{frame}{Derivadas}
  \begin{align*}
    \big(\sin(x)\big)' & = \cos(x)        \\[2mm]
    \big(\cos(x)\big)' & = -\sin(x)       \\[2mm]
    \big(\tan(x)\big)' & = \sec^2(x)      \\[2mm]
    \big(\sec(x)\big)' & = \tan(x)\sec(x) \\[2mm]
    \big(\csc(x)\big)' & = \cot(x)\csc(x) \\[2mm]
    \big(\cot(x)\big)' & = \csc^2(x)
  \end{align*}
\end{frame}

%------------------------------------------------------------------------------%
\section{Senos e Cossenos}

\begin{frame}{Senos $\times$ cossenos}
  Estratégias para calcular
  \[
    \int \sin^m(x) \cos^n(x)\,dx
  \]

  Depende se $m$ ou $n$ são pares ou ímpares
\end{frame}

%------------------------------------------------------------------------------%
\begin{frame}{Potência do $\sin(x)$ é ímpar}
  \onslide<+->{Potência do $\sin(x)$ é ímpar}
  \begin{enumerate}[<+->]
    \setlength\itemsep{1em}
    \item Guarde um fator $\sin(x)$
    \item Use que $\sin^2(x) = 1 - \cos^2(x)$
    \item Faça a substituição $u=\cos(x)$
  \end{enumerate}
\end{frame}

%------------------------------------------------------------------------------%
\begin{frame}{Potência do $\cos(x)$ é ímpar}
  \onslide<+->{Potência do $\cos(x)$ é ímpar}
  \begin{enumerate}[<+->]
    \setlength\itemsep{1em}
    \item Guarde um fator $\cos(x)$
    \item Use que $\cos^2(x) = 1 - \sin^2(x)$
    \item Faça a substituição $u=\sin(x)$
  \end{enumerate}
\end{frame}

%------------------------------------------------------------------------------%
\begin{frame}{Ambas são pares}
  \onslide<+->{Ambas são pares}

  \onslide<+->{Use
  \[
    \sin^2(x)= \frac{1}{2}\big(1-\cos(2x)\big) \qquad\text{ou}\qquad
    \cos^2(x)= \frac{1}{2}\big(1+\cos(2x)\big)
  \]}

  \onslide<+->{Transformação útil
  \[
    \sin(2x) = 2\sin(x)\cos(x)
  \]}
\end{frame}

% 1
%------------------------------------------------------------------------------%
\addtocounter{exemplo}{1}
\section{Exemplo \theexemplo}

%------------------------------------------------------------------------------%
\begin{frame}{Exemplo \theexemplo}
  \onslide<+->{Calcule \(\int \sin^3(x)\cos^2(x)\,dx\)}

  \vspace{\baselineskip}
  \onslide<+->{Integral do tipo
  \[
    \int \sin^m(x) \cos^n(x)\,dx \qquad m \;\text{ ímpar}
  \]}

  \onslide<+->{Guardar $\sin(x)$}

  \onslide<+->{Usar \(\sin^2(x)= 1- \cos^2(x)\)}
\end{frame}

%------------------------------------------------------------------------------%
\begin{frame}{Exemplo \theexemplo}
  \begin{align*}
    \onslide<+->{F & = \int \sin^3(x) \cos^2(x) \,dx                            \\[2mm]}
    \onslide<+->{  & = \int \sin^2(x)                 \cos^2(x) \sin(x) \,dx    \\[2mm]}
    \onslide<+->{  & = \int \left( 1 - \cos^2(x) \right) \cos^2(x) \sin(x) \,dx}
  \end{align*}
  \onslide<+->{Substituição simples
  \[
    u  =  \cos(x)     \qquad
    du = -\sin(x)dx
  \]}
  \onslide<+->{\[
    F = \int \left( 1 - u^2 \right) u^2 \,(-du)
  \]}
\end{frame}

%------------------------------------------------------------------------------%
\begin{frame}{Exemplo \theexemplo}
  \begin{align*}
    \onslide<+->{F & = \int \left( 1 - u^2 \right) u^2 \,(-du) \\[2mm]}
    \onslide<+->{  & = \int \left( u^2 - 1\right) u^2 \,du     \\[2mm]}
    \onslide<+->{  & = \int u^4 - u^2 \, du                    \\[2mm]}
    \onslide<+->{  & = \frac{u^5}{5} - \frac{u^3}{3} + C       \\[2mm]}
    \onslide<+->{  & = \frac{1}{5}\cos^5(x) - \frac{1}{3}\cos^3(x) + C}
  \end{align*}
\end{frame}

% 2
%------------------------------------------------------------------------------%
\addtocounter{exemplo}{1}
\section{Exemplo \theexemplo}

%------------------------------------------------------------------------------%
\begin{frame}{Exemplo \theexemplo}
  \onslide<+->{Encontre \(\int \cos^3(x)\,dx\)}

  \vspace{\baselineskip}
  \onslide<+->{Integral do tipo
  \[
    \int \sin^m(x) \cos^n(x)\,dx \qquad m=0 \;\text{ e }\; n \;\text{ ímpar}
  \]}

  \onslide<+->{Guardar $\cos(x)$}

  \onslide<+->{Usar \(\cos^2(x)= 1- \sin^2(x)\)}
\end{frame}

%------------------------------------------------------------------------------%
\begin{frame}{Exemplo \theexemplo}
  \begin{align*}
    \onslide<+->{F & = \int \cos^3(x)\,dx                            \\[2mm]}
    \onslide<+->{  & = \int \cos^2(x) \cos(x)\,dx                    \\[2mm]}
    \onslide<+->{  & = \int \left( 1 - \sin^2(x) \right) \cos(x)\,dx}
  \end{align*}
  \onslide<+->{Substituição simples
  \[
    u  = \sin(x)     \qquad
    du = \cos(x)dx
  \]}
  \onslide<+->{\[
    F = \int \left( 1 - u^2 \right) \,du
  \]}
\end{frame}

%------------------------------------------------------------------------------%
\begin{frame}{Exemplo \theexemplo}
  \begin{align*}
    \onslide<+->{F & = \int 1 - u^2 \,du                  \\[2mm]}
    \onslide<+->{  & = u - \frac{u^3}{3} + C              \\[2mm]}
    \onslide<+->{  & = \sin(x) - \frac{1}{3}\sin^3(x) + C}
  \end{align*}
\end{frame}

% 3
%------------------------------------------------------------------------------%
\addtocounter{exemplo}{1}
\section{Exemplo \theexemplo}

%------------------------------------------------------------------------------%
\begin{frame}{Exemplo \theexemplo}
  \onslide<+->{Calcule \(\int_0^{\pi} \sin^2(x) dx\)}

  \vspace{\baselineskip}
  \onslide<+->{Integral do tipo
  \[
    \int \sin^m(x) \cos^n(x)\,dx \qquad n=0 \;\text{ e }\; m \;\text{ par}
  \]}

  \onslide<+->{Usar \(\sin^2(x)= \frac{1}{2}\big( 1 - \cos(2x) \big)\)}
\end{frame}

%------------------------------------------------------------------------------%
\begin{frame}{Exemplo \theexemplo}
  \begin{align*}
    \onslide<+->{\int_0^{\pi} \sin^2(x)\, dx
    & = \int_0^\pi \frac{1}{2}\big( 1 - \cos(2x) \big) dx                 \\[2mm]}
    \onslide<+->{& = \frac{1}{2} \int_0^\pi 1-\cos(2x)\, dx                            \\[2mm]}
    \onslide<+->{& = \frac{1}{2} \left( x - \frac{\sin(2x)}{2} \right) \evalat{0}{\pi} \\[2mm]}
    \onslide<+->{& = \frac{1}{2}
        \left[
           \left( \pi - \frac{\sin(2\pi)}{2} \right)
          -\left(   0 - \frac{\sin(0)}{2}    \right)
        \right]                                                           \\[2mm]}
    \onslide<+->{& = \frac{\pi}{2}}
  \end{align*}
\end{frame}

% 4
%------------------------------------------------------------------------------%
\addtocounter{exemplo}{1}
\section{Exemplo \theexemplo}

%------------------------------------------------------------------------------%
\begin{frame}{Exemplo \theexemplo}
  \onslide<+->{Calcule \(\int \sin^2(x) \cos^2(x) dx\)}

  \vspace{\baselineskip}
  \onslide<+->{Integral do tipo
  \[
    \int \sin^m(x) \cos^n(x)\,dx \qquad m \;\text{ e }\; n \;\text{ pares}
  \]}

  \onslide<+->{Usar
  \begin{gather*}
    \sin^2(x) = \frac{1}{2}\big(1-\cos(2x)\big) \\[2mm]
    \cos^2(x) = \frac{1}{2}\big(1+\cos(2x)\big)
  \end{gather*}}
\end{frame}

%------------------------------------------------------------------------------%
\begin{frame}{Exemplo \theexemplo}
  \begin{align*}
    \onslide<+->{F & = \int \sin^2(x) \cos^2(x) dx \\[2mm]}
    \onslide<+->{  & = \int \frac{1}{2} \big( 1 - \cos(2x) \big)
                            \frac{1}{2} \big( 1 + \cos(2x) \big) dx   \\[2mm]}
    \onslide<+->{  & = \frac{1}{4} \int 1 - \cos^2(2x)\, dx & & (a-b)(a+b) = a^2 - b^2 \\[2mm]}
    \onslide<+->{  & = \frac{1}{4} \int \sin^2(2x) dx       & & \sin^2(\alpha) + \cos^2(\alpha)=1\\[2mm]}
    \onslide<+->{  & = \frac{1}{4} \int \frac{1}{2} \big( 1 - \cos(4x) \big) dx }
  \end{align*}
\end{frame}

%------------------------------------------------------------------------------%
\begin{frame}{Exemplo \theexemplo}
\begin{align*}
  \onslide<+->{F & = \frac{1}{4} \int \frac{1}{2} \big( 1 - \cos(4x) \big) dx \\[2mm]}
  \onslide<+->{  & = \frac{1}{8} \left(\int 1 dx - \int \cos(4x) dx\right)  \\[2mm]}
  \onslide<+->{  & = \frac{1}{8} \left(x - \frac{\sin(4x)}{4} + c \right)   \\[2mm]}
  \onslide<+->{  & = \frac{x}{8} - \frac{\sin(4x)}{32} + c}
\end{align*}
\end{frame}

%%------------------------------------------------------------------------------%
%\begin{frame}{Exemplo \theexemplo}
%  \begin{align*}
%    \onslide<+->{G & = \int \cos^2(2x) dx                                     \\[2mm]}
%    \onslide<+->{  & = \int \frac{1}{2} \big( 1 + \cos(4x) \big)dx
%      & \cos^2(\alpha) = \frac{1}{2}\big( 1+ \cos(2\alpha) \big) \\[2mm]}
%    \onslide<+->{  & = \int \frac{1}{2} - \frac{1}{2} \cos(4x) \,dx           \\[2mm]}
%    \onslide<+->{  & = \frac{x}{2} - \frac{1}{8} \sin(4x) + C}
%  \end{align*}
%\end{frame}
%
%%------------------------------------------------------------------------------%
%\begin{frame}{Exemplo \theexemplo}
%  \begin{align*}
%    \onslide<+->{F & = \frac{x}{4} - \frac{1}{4}\int \cos^2(2x)\, dx                                 \\[2mm]}
%    \onslide<+->{  & = \frac{x}{4} - \frac{G}{4}                                                     \\[2mm]}
%    \onslide<+->{  & = \frac{x}{4} - \frac{1}{4} \left(\frac{x}{2} - \frac{1}{8} \sin(4x) + C\right) \\[2mm]}
%    \onslide<+->{  & = \frac{x}{4} - \frac{x}{8} + \frac{1}{64} \sin(4x) + C_1                       \\[2mm]}
%    \onslide<+->{  & = \frac{x}{4} + \frac{1}{64} \sin(4x) + C_1}
%  \end{align*}
%\end{frame}

%------------------------------------------------------------------------------%
\section{Lista Mínima}

%------------------------------------------------------------------------------%
\begin{frame}{Lista Mínima}
  Estudar a Seção 4.3 da Apostila

  Exercícios: 1a-f, 2, 3, 5, 6

  \vfill
  Atenção: A prova é baseada no livro, não nas apresentações
\end{frame}

%------------------------------------------------------------------------------%
\end{document}
%------------------------------------------------------------------------------%
